\section{Stable coordinates and a physical defense against guide collapse}
\label{sec:robust-guide}

The purpose of the Gaussian guide is to make a conditional TT proposal easier
to fit. A failed guide should therefore make approximation less efficient,
not destroy the coordinates in which the proposal is represented. This
distinction matters for the four-dimensional stochastic-volatility example.
In the A09 confirmation sequence s04 at time 8, a nine-node signed sparse
quadrature rule assigned normalized weight essentially one to a single node.
The sum of the other absolute weights was $1.38\,10^{-31}$. Its computed
covariance eigenvalues lay between $3.72\,10^{-32}$ and $2.44\,10^{-31}$,
although the predictive covariance eigenvalues were between 1.20 and 1.32.
All higher tested rules had an indefinite covariance. A check relative to
the candidate covariance's own scale accepted the tiny positive matrix.
Consequently the Gaussian component used to defend the TT polynomial had
the same collapsed physical scale as the fitted TT. The corrected filter
then failed log-evidence agreement, despite apparently valid factorizations.

The highest-resolution particle reference for this sequence had marginal
variances approximately $(1.07,0.25,0.71,0.90)$ at that time. This is strong
diagnostic evidence against interpreting the near-zero guide covariance as
real posterior concentration. It is not a theorem about every small
posterior covariance. The saved replay is
\path{docs/plans/artifacts/observation-tt-warm-improvement-20260916-01/covariance-collapse-diagnostic.md}.
The A09 evidence also exposed seed reuse in earlier attempts; earlier
independence-based Monte Carlo standard errors must not be reused without
their own seed audit. The present construction and its fresh tests are
specified by amendment A10.

We use three protections with different mathematical purposes. First,
quadrature diagnostics reject numerically unresolved Gaussian moments.
Second, model-based TT coordinates remain nonsingular even when those
moments fail. Third, a mixture with the physical transition density gives
the proposal independent tail coverage. None of these statements alone
certifies the accuracy of Gaussian quadrature or the quality of TT fitting.
They define an optional extension of the proposal construction, rather than
a new claim about fidelity to the Zhao--Cui author implementation.

\subsection{What a covariance check can establish}

Let the Gaussian predictive approximation be $p^-(x)=\mathcal N(x;m,S)$,
where $S=LL^\top\succ0$. A quadrature rule has nodes $u_i$ in predictive
standard coordinates and possibly signed weights $w_i$. Put
\begin{equation}
 a_i=w_i\exp\{\ell(m+Lu_i)-b\},\qquad
 A=\sum_i a_i,\qquad b=\max_i\ell(m+Lu_i),\qquad
 \alpha_i=a_i/A .
 \label{eq:robust-signed-weights}
\end{equation}
Here $\ell(x)=\log g(y\mid x)$ and $b$ only prevents numerical overflow.
When $A>0$, the formal quadrature mean and covariance in these coordinates
are $\bar u=\sum_i\alpha_i u_i$ and
$C=\sum_i\alpha_i(u_i-\bar u)(u_i-\bar u)^\top$.
The physical moments are $m+L\bar u$ and $LCL^\top$. These formulas are
quadrature estimates; signed $\alpha_i$ are not probabilities.

\begin{proposition}[A relative SPD test does not protect physical scale]
\label{prop:robust-scale}
For any $c\in(0,1)$, the test
$\lambda_{\min}(C)>c\|C\|_2$ accepts every $C=\delta I_d$ with
$\delta>0$, however small $\delta$ is relative to a predictive covariance.
Moreover, positive normalized quadrature weights do not prevent their
weighted covariance from converging to zero.
\end{proposition}
\begin{proof}
For $C=\delta I_d$, both $\lambda_{\min}(C)$ and $\|C\|_2$ equal
$\delta$, so the inequality reduces to $1>c$. For the second assertion,
fix finitely many nodes and let the weight of one node tend to one while
all other nonnegative weights tend to zero. The weighted mean tends to that
node, and every term in the finite covariance sum tends to zero.
Thus positivity prevents an indefinite exact weighted covariance but does
not prevent concentration on a single quadrature node.
\end{proof}

There are two distinct numerical warnings. Cancellation is measured by
$\kappa_A=\sum_i|a_i|/|A|$; dominance is measured by
$D_A=\max_i|a_i|/\sum_i|a_i|$. The latter uses absolute weights and
remains meaningful for a signed rule. It is not an effective sample size.

\begin{proposition}[Cancellation and moment sensitivity]
\label{prop:robust-cancellation}
Suppose the summands $a_i$ are exact, their floating-point sum is computed
without overflow or underflow, and $n\epsilon_{\rm mach}<1$. In the
usual rounding model its absolute error is at most
$\gamma_n\sum_i|a_i|$, where
$\gamma_n=n\epsilon_{\rm mach}/(1-n\epsilon_{\rm mach})$.
Hence the relative error bound is $\gamma_n\kappa_A$ when $A\ne0$.
In exact arithmetic a signed normalized weighted covariance can be
indefinite even when $A>0$; this is a separate failure from cancellation.
\end{proposition}
\begin{proof}
Expanding successive rounded additions gives
$\widehat A=\sum_i a_i(1+\theta_i)$ with
$|\theta_i|\le\gamma_n$. The triangle inequality proves the absolute
bound, and division by $|A|$ proves the relative bound. The same argument
applies entry by entry to sums defining moments, with the appropriate
absolute summands. It does not account for errors in evaluating likelihoods
or centered nodes. For indefiniteness, the one-dimensional nodes $0,1$
with normalized weights $2,-1$ have mean $-1$ and variance
$2(0+1)^2-(1+1)^2=-2$. Their total weight is one.
\end{proof}

We therefore evaluate covariance in predictive units and compare its
smallest eigenvalue with a margin that has an independent unit scale:
\begin{equation}
 \eta_C=64\epsilon_{\rm mach} n
 \max\!\left\{1,\sum_i |\alpha_i|\|u_i-\bar u\|^2\right\},
 \qquad \lambda_{\min}(C)>\eta_C .
 \label{eq:robust-margin}
\end{equation}
Finite values, positive signed mass and a corresponding cancellation check
are also required. The factor 64 is a conservative numerical hypothesis,
not a proved forward-error constant for the whole algorithm. This check
rejects numerically unusable guide moments; it does not assert that a true
posterior cannot be concentrated. No eigenvalue is clipped and no hidden
ridge is added to accepted moments. Dominance, cancellation, eigenvalues
and standardized third and fourth moments are retained as diagnostics.

Numerical admissibility is followed by a resolution check. Successive
admissible rules must agree in mean, covariance and log integral to declared
tolerances: $0.02$ predictive standard-deviation units for the Euclidean
mean difference, $0.05$ for the relative Frobenius covariance difference,
and $0.02$ nats for the log integral. The covariance denominator is the
larger of one and the earlier covariance's Frobenius norm in predictive
units. Agreement of two rules is evidence of resolution at those orders,
not a quadrature error bound: both rules can miss the same mass. A failed
check triggers the following independent construction.

\subsection{An observation-adapted positive quadrature repair}

In the example, $x\in\mathbb R^d$, $\beta>0$, and, conditionally on $x$,
$y_j$ are independent $\mathcal N(0,\beta^2 e^{x_j})$ variables. Combining
this likelihood with the Gaussian predictive approximation gives the
negative log density, up to a constant independent of $x$,
\begin{equation}
 U(x)=\tfrac12(x-m)^\top S^{-1}(x-m)
       +\tfrac12\sum_{j=1}^d\{x_j+c_j e^{-x_j}\},
 \qquad c_j=(y_j/\beta)^2\ge0 .
 \label{eq:robust-sv-potential}
\end{equation}
This is a density for the Gaussian-predictive approximation. It equals
neither the full filtering density in general nor a TT-retained target.

\begin{proposition}[Unique mode of the Gaussian-predictive SV update]
\label{prop:robust-mode}
If $S\succ0$, $\beta>0$ and all $y_j$ are finite, $U$ is coercive and
strongly convex. It has a unique minimizer $x_*$, with
\begin{align}
 \nabla U(x)&=S^{-1}(x-m)+\tfrac12(\boldsymbol1-c\odot e^{-x}),
 \nonumber\\
 H(x)&=S^{-1}+\tfrac12\operatorname{diag}(c\odot e^{-x})\succ0 .
 \label{eq:robust-mode-derivatives}
\end{align}
The mode is a smooth function locally of $(m,S,y,\beta)$ on the domain
$S\succ0$, $\beta>0$, including observations with zero components.
\end{proposition}
\begin{proof}
Differentiating \eqref{eq:robust-sv-potential} gives the displayed gradient
and Hessian. The diagonal addition is positive semidefinite, so
$H(x)\succeq S^{-1}\succeq\lambda_{\min}(S^{-1})I$ for all $x$.
The quadratic term dominates the possibly negative linear term as
$\|x\|\to\infty$, while the exponential term is nonnegative. This
proves coercivity, existence and uniqueness. The gradient is smooth in all
parameters on the stated open domain. Its derivative with respect to $x$
at the solution is the invertible Hessian. The implicit function theorem
gives the local smooth dependence. Setting a $y_j$ to zero leaves these
arguments unchanged.
\end{proof}

We solve for the mode with damped Newton steps and a backtracking decrease
test using the displayed analytical derivatives. A finite residual check
is required before using $V=H(x_*)^{-1}$. The proposition concerns the
exact mode; it does not prove that any fixed numerical iteration budget
will attain its tolerance. Failed factorizations or unresolved residuals
are recorded rather than concealed by covariance clipping.

Let $q_L(x)=\mathcal N(x;x_*,V)$ and write $x=x_*+Bu$, where
$BB^\top=V$ and $u$ has standard Gaussian density $\rho$.

\begin{proposition}[Positive quadrature with the correct density ratio]
\label{prop:robust-change-measure}
For any function $h$ integrable under $p^-(x)g(y\mid x)$,
\begin{equation}
 \int h(x)p^-(x)g(y\mid x)\,dx
 =\int h(x_*+Bu)
   \frac{p^-(x_*+Bu)g(y\mid x_*+Bu)}{q_L(x_*+Bu)}\rho(u)\,du .
 \label{eq:robust-change-measure}
\end{equation}
A positive Gauss--Hermite rule for the right-hand side produces positive
normalized weights. With strictly positive finite likelihood ratios and
nodes whose physical locations affinely span $\mathbb R^d$, its exact
weighted covariance is positive definite. Neither assertion establishes
the accuracy or numerical conditioning of the finite quadrature rule.
\end{proposition}
\begin{proof}
The substitution has Jacobian $|\det B|$, and
$q_L(x_*+Bu)|\det B|=\rho(u)$. Substitution proves the identity and
shows why the predictive-to-Laplace density ratio cannot be omitted.
Multiplying positive quadrature weights by a positive likelihood ratio
preserves positivity. For normalized weights $\omega_i>0$ and any
$v\ne0$, the weighted variance is
$\sum_i\omega_i[v^\top(x_i-\bar x)]^2$. It vanishes only if
$v^\top x_i$ is constant at every node. Affine spanning excludes this
possibility. Proposition~\ref{prop:robust-scale} still permits arbitrarily
small variances when almost all weight is concentrated at one node.
\end{proof}

A10 tests tensor orders 5, 7 and 9, at most $9^4=6561$ nodes, using
log-sum-exp normalization and the same numerical and resolution diagnostics.
The small-dimensional tensor rule is a bounded repair experiment; its
exponential growth prevents a claim of a scalable high-dimensional
implementation. It is not the generic all-axes retained-grid Zhao--Cui
route. If neither signed sparse quadrature nor this repair resolves the
guide, the predictive Gaussian is used explicitly as an unresolved-guide
fallback. This preserves a usable proposal but supplies no observation
adaptation at that step.

\subsection{TT coordinates that do not inherit guide covariance collapse}

Let the physical dynamics be $x_t=A x_{t-1}+\xi_t$, with
$\xi_t\sim\mathcal N(0,Q)$, $Q\succ0$, and initial covariance
$P_0\succ0$. Define deterministic model covariances
\begin{equation}
 R_0=P_0,\qquad R_t=A R_{t-1}A^\top+Q,\qquad
 x_t=c_t+L_t u_t,\quad L_tL_t^\top=R_t .
 \label{eq:robust-chart}
\end{equation}
The center $c_t$ may be the accepted or repaired guide mean. The covariance
$R_t$ is independent of the guide's observation-weighted covariance. In the
stationary fixture, $A P_0A^\top+Q=P_0$ and $R_t=P_0$ for every $t$.

\begin{proposition}[A nonsingular chart and invariant physical proposals]
\label{prop:robust-chart}
For finite $t$, finite $A,Q,P_0$ satisfying the assumptions above,
$R_t\succ0$ and $R_t\succeq Q$ for $t\ge1$.
The affine change of coordinates is bijective. If a physical density $p$
is represented in these coordinates as
$p_U(u)=p(c_t+L_tu)|\det L_t|$, then its physical law is unchanged.
The same statement holds for a joint law after transforming both states.
\end{proposition}
\begin{proof}
For nonzero $v$,
$v^\top R_tv=(A^\top v)^\top R_{t-1}(A^\top v)+v^\top Q>0$;
induction starts from $P_0\succ0$. The first summand is nonnegative, which
also proves $R_t\succeq Q$. Every finite recursion has finite entries in
exact real arithmetic, and its Cholesky factor is nonsingular. The
change-of-variables formula gives the asserted density, with determinant
$|\det L_t|$; a block-diagonal affine transformation gives the joint
version. Centering translates the coordinates and does not change the
factor's singular values.
\end{proof}

This proposition removes guide collapse as a source of chart singularity.
It does not bound $R_t$ uniformly for unstable dynamics, guarantee safe
floating-point arithmetic for arbitrary $A,Q$, or imply that a low-degree
TT fits well in the broader coordinates. Numerical factorizations are still
checked. In particular, Gaussian initialization and regression sampling must
transform the same physical SGQF joint into the new charts. Reinterpreting
$R_t$ as the guide covariance would instead change the initializer and
the design distribution. The implementation checks this distinction at the
consumer endpoint, not just inside a coordinate utility.

For clarity, let the incoming guide be $\mathcal N(m_z,P_z)$, the current
guide be $\mathcal N(m_x,P_x)$, and define
$S_z=A P_zA^\top+Q$, $K=P_zA^\top S_z^{-1}$ and
$B_z=P_z-KS_zK^\top$. The Gaussian joint used for initialization has
\begin{equation}
 \begin{pmatrix}x\\z\end{pmatrix}\sim\mathcal N\!\left(
 \begin{pmatrix}m_x\\m_z+K(m_x-Am_z)\end{pmatrix},
 \begin{pmatrix}P_x&P_xK^\top\\KP_x&B_z+KP_xK^\top\end{pmatrix}
 \right).
 \label{eq:robust-initializer-joint}
\end{equation}
This follows by drawing $x$ from the current guide and then drawing
$z\mid x\sim\mathcal N(m_z+K(x-Am_z),B_z)$; conditional expectation
and the law of total covariance give every block. Equation
\eqref{eq:robust-chart} transforms this joint without substituting $R_t$
for either $P_x$ or $P_z$.

\subsection{A defensive component in physical coordinates}

Let $q_{\rm TT}(x\mid z)$ be the normalized conditional proposal obtained
from the fitted pair TT, including its internal polynomial-zero defense.
That internal component is Gaussian in the TT chart. Define instead the
additional physical mixture
\begin{equation}
 q_\varepsilon(x\mid z)
 =(1-\varepsilon)q_{\rm TT}(x\mid z)
       +\varepsilon f(x\mid z),\qquad 0<\varepsilon<1,
 \quad f(x\mid z)=\mathcal N(x;Az,Q).
 \label{eq:robust-mixture}
\end{equation}
At time zero the second component is the physical prior. A component label
is sampled with probabilities $1-\varepsilon,\varepsilon$, followed by a
draw from that component. Both component densities must then be evaluated
at the actual sampled point. Using only the sampled component's density in
the denominator implements a different importance calculation.

\begin{proposition}[Normalization and exact conditional importance correction]
\label{prop:robust-mixture-correction}
Suppose $q_{\rm TT}(\cdot\mid z)$ is a probability density and
$Z_y(z)=\int f(x\mid z)g(y\mid x)\,dx\in(0,\infty)$.
Then \eqref{eq:robust-mixture} is a probability density with
$q_\varepsilon\ge\varepsilon f$. For every integrable $h$,
\begin{equation}
 \mathbb E_{q_\varepsilon}\!\left[
 h(X)\frac{f(X\mid z)g(y\mid X)}{q_\varepsilon(X\mid z)}\right]
 =\int h(x)f(x\mid z)g(y\mid x)\,dx .
 \label{eq:robust-importance}
\end{equation}
Thus the exact importance target is unchanged by this proposal mixture.
\end{proposition}
\begin{proof}
Integration of the convex combination gives one, and nonnegativity gives
the lower bound. Multiplying the integrand on the left by its sampling
density cancels the denominator. Since $f$ is everywhere positive, the
mixture covers the target support. Integrability permits the expectation.
\end{proof}

The equality is for unnormalized importance integrals conditional on $z$.
It does not make a finite-sample self-normalized mean unbiased, nor does it
validate a TT approximation's own normalizing constant. In the particle
filter the usual previous particle weights and the declared resampling law
remain part of the calculation.

\begin{proposition}[Finite conditional second moment for the SV likelihood]
\label{prop:robust-second-moment}
Under the SV likelihood above, finite $z,y$, $\beta>0$ and $Q\succ0$,
the importance weight $W=f(X\mid z)g(y\mid X)/q_\varepsilon(X\mid z)$
has finite conditional second moment. In particular, with $\mu=Az$,
\begin{equation}
 \mathbb E_{q_\varepsilon}[W^2\mid z]
 \le \frac{1}{\varepsilon}\mathbb E_{f(\cdot\mid z)}[g(y\mid X)^2]
 \le \frac{(2\pi\beta^2)^{-d}}{\varepsilon}
 \exp\!\left\{-\boldsymbol1^\top\mu
                  +\tfrac12\boldsymbol1^\top Q\boldsymbol1\right\}.
 \label{eq:robust-second-moment}
\end{equation}
The normalized conditional importance ratio $W/Z_y(z)$ also has finite
second moment. This is a pointwise-in-$(y,z)$ assertion, not a uniform ESS
bound or a lower bound on any realized finite-particle ESS.
\end{proposition}
\begin{proof}
Since $q_\varepsilon\ge\varepsilon f$,
$\int f^2g^2/q_\varepsilon\le\varepsilon^{-1}\int fg^2$.
Squaring the likelihood gives
\[
 g(y\mid x)^2=(2\pi\beta^2)^{-d}
 \exp\!\left\{-\boldsymbol1^\top x-\sum_j c_j e^{-x_j}\right\}
 \le (2\pi\beta^2)^{-d}e^{-\boldsymbol1^\top x}.
\]
The Gaussian moment-generating function at $-\boldsymbol1$ gives the
second bound. The same bound, or Cauchy--Schwarz, makes $Z_y(z)$ finite;
strict positivity of the likelihood and transition makes it positive.
Division by its square proves the normalized assertion. It may be very
small, so this argument supplies no useful universal numerical bound on
the normalized variance. At time zero the proof is identical with the
prior mean and covariance replacing $\mu,Q$.
\end{proof}

The first inequality bounds the second moment by $1/\varepsilon$ times
that of the transition proposal. A10 evaluates $\varepsilon=0.05,0.2,0.5$,
whose multipliers are 20, 5 and 2. These are declared robustness/cost
hypotheses. They do not predict which setting has the smallest filtering
error. The smallest setting passing the calibration safety/non-harm checks
is frozen before fresh confirmation. Runtime switching based on heldout
performance is excluded.

\subsection{Differentiability and the remaining empirical question}

\begin{proposition}[Smooth mixture density on a fixed branch]
\label{prop:robust-smoothness}
Let $\vartheta$ collect continuous state or model parameters. On an open
domain suppose $q_{\rm TT}(x\mid z;\vartheta)$ and
$f(x\mid z;\vartheta)$ are positive $C^1$ densities and $\varepsilon$
is a fixed number in $(0,1)$. Then $\log q_\varepsilon$ is $C^1$, with
\begin{equation}
 \nabla_\vartheta\log q_\varepsilon
 =r_{\rm TT}\nabla_\vartheta\log q_{\rm TT}
   +r_f\nabla_\vartheta\log f,\qquad
 r_{\rm TT}=\frac{(1-\varepsilon)q_{\rm TT}}{q_\varepsilon},
 \quad r_f=\frac{\varepsilon f}{q_\varepsilon}.
 \label{eq:robust-mixture-score}
\end{equation}
This statement alone gives no derivative of adaptive guide selection,
TT fitting, component selection or particle resampling.
\end{proposition}
\begin{proof}
A fixed positive convex combination of $C^1$ functions is positive and
$C^1$. Differentiating its logarithm and multiplying and dividing each
summand by its positive component density gives the formula. The argument
treats the component densities as differentiable functions on the stated
domain; it does not differentiate any algorithm used to select or construct
them. Branch changes, active-set changes in L1 fitting, singular-value
choices and discrete resampling therefore require additional analysis.
\end{proof}

The proposal remains analytically differentiable in its evaluation
coordinates when its fitted cores and positive charts are held fixed. A
total analytical derivative of the complete finite filtering program is a
separate, unproved requirement. The immediate empirical question is more
limited: do the numerical guards detect the observed failure, do the new
charts and physical defense preserve valid filtering on fresh sequences,
and is the cost acceptable? Tests retain quadrature failures, normalized
filtering MSE, log-evidence agreement, particle ESS, maximum weights,
ancestry and timing. Loss against a Gaussian joint describes approximation
quality; it cannot by itself decide whether a TT proposal is useful after
exact importance correction.

